% !TeX root = ../../Book3.tex
\section{Riemann 映照定理}
\label{b3:sec:0503}

前节的极值构造给出存在性。唯一性则来自两个映射之间的圆盘自同构：把一个映射接在另一个的逆后面，就把区域的形状全部消去，只留下归一化能否固定圆盘坐标的问题。

\subsection{存在性}
\label{b3:subsec:050301}

\term[riemann-yingzhao-dingli]{Riemann 映照定理}[Riemann mapping theorem]在本卷采用真平面区域版本。
\begin{theorem}[Riemann 映照]\label{b3:thm:05-riemann}
设 \(U\subsetneq\mathbb C\) 是单连通区域，\(a\in U\)。则存在唯一双全纯映射 \(f:U\rightarrow\mathbb D\)，满足 \(f(a)=0\)、\(f'(a)>0\)。
\end{theorem}
\begin{proof}
存在性由归一化族非空、极值可达与满射性三引理得到。单叶函数导数非零，所以局部逆全纯；这些局部逆在重叠处因单射一致，拼成全域全纯逆。唯一性见紧随其后的论证。
\end{proof}
存在定理不要求边界光滑，因此即使边界极粗糙，内部共形坐标仍存在；它不自动延拓为欧氏闭包同胚。边界延拓将在\chapref{b3:ch:07}另列条件。

\begin{proposition}\label{b3:prop:05-exhaustion-maps}
设 \(U_1\subset U_2\subset\cdots\) 为真单连通区域，\(a\in U_1\)，且 \(U=\bigcup_n U_n\subsetneq\mathbb C\)。令 \(f_n:U_n\rightarrow\mathbb D\) 和 \(f:U\rightarrow\mathbb D\) 在 \(a\) 按本章归一化。则 \(f_n\) 在 \(U\) 的每个紧集上最终有定义并一致趋于 \(f\)，且 \(r_{U_n}(a)\uparrow r_U(a)\)。
\end{proposition}
\begin{proof}
每个紧集包含于某个 \(U_N\)，于是局部有界与对角抽取给 \(f_n\) 的一个全纯极限 \(g\)。导数单调且不小于 \(f'(a)>0\)，故极限不常值，由单叶极限定理仍单叶并映入圆盘。令 \(h_n=f_n^{-1}\)。函数 \(k_n=f\circ h_n\) 是圆盘自映射且固定0，所以 \(|k_n(z)|\le|z|\)。因此对固定紧盘 \(|z|\le r\)，\(h_n(z)\) 落在 \(f^{-1}(\overline{D(0,r)})\) 内，后者是 \(U\) 内的紧集。对 \(k_n\) 用 Montel，再复合 \(f^{-1}\)，得到 \(h_n\) 的紧极限 \(h:\mathbb D\rightarrow U\)。其导数在0不为零，因为 \(f_n'(a)\) 有正的有限上界。沿同一个子列取极限，在固定紧集上由 \(f_n\circ h_n=\mathrm{id}\) 和 \(h_n\circ f_n=\mathrm{id}\) 得 \(g\circ h=\mathrm{id}\)、\(h\circ g=\mathrm{id}\)。所以 \(g\) 是归一化双全纯映射，唯一性给 \(g=f\)。任何子列都有同一个极限，故全列收敛，导数与共形半径也随之收敛。
\end{proof}
这个命题在逐步逼近复杂区域时很有用，但它只断言内部紧集上的收敛。区域边界可能在逼近过程中改变，映射的边界值不能由此直接交换极限。连通性和单连通性对每个 \(U_n\) 都有要求；任意截取 \(U\cap D(0,n)\) 可能不连通，因此不能未经检查便作为符合假设的穷竭。

\subsection{唯一性}
\label{b3:subsec:050302}

\begin{proof}[\cref{b3:thm:05-riemann}的唯一性证明]
若 \(f,g\) 都满足条件，\(h=g\circ f^{-1}\) 为圆盘自同构且 \(h(0)=0\)。对 \(h\) 和 \(h^{-1}\) 分别用 Schwarz 引理，得到 \(|h'(0)|\le1\) 与 \(|(h^{-1})'(0)|\le1\)。两导数相乘为1，故模都等于1，Schwarz 等号情形给 \(h(z)=cz\)、\(|c|=1\)。另一方面 \(h'(0)=g'(a)/f'(a)>0\)，所以 \(c=1\)，得到 \(g=f\)。
\end{proof}
若只固定 \(f(a)=0\)，仍有一个单位复数旋转自由度；若连 \(a\) 的像都不固定，则还允许所有圆盘自同构。归一化的每一条件分别消除不同自由度，不能把“唯一”理解为所有共形映射完全相同。

\begin{definition}\label{b3:def:05-conformal-radius}
设 \(U\subsetneq\mathbb C\) 单连通，\(a\in U\)，\(F:\mathbb D\rightarrow U\) 双全纯且 \(F(0)=a\)。定义 \(U\) 在 \(a\) 的\term[gongxing-ban-jing]{共形半径}[conformal radius]为 \(r_U(a)=|F'(0)|\)。
\end{definition}
两个这样的逆映射相差圆盘旋转，所以定义与选择无关。若 \(f\) 是本章归一化映射，则 \(r_U(a)=1/f'(a)\)。对圆盘 \(D(c,R)\)，将 \((z-c)/R\) 与圆盘自同构复合得到
\[
 r_{D(c,R)}(a)=\frac{R^2-|a-c|^2}{R}.
\]
对右半平面，归一化坐标可取 \((z-a)/(z+\overline a)\) 再校正相位，故 \(r_U(a)=2\operatorname{Re}a\)。它是一个随点变化的内在尺度，并非到边界的欧氏距离。

\begin{proposition}\label{b3:prop:05-radius-monotonicity}
若 \(U\subset V\subsetneq\mathbb C\) 为单连通区域且 \(a\in U\)，则 \(r_U(a)\le r_V(a)\)，真包含时严格。双全纯 \(T:U\rightarrow W\) 满足 \(r_W(Ta)=|T'(a)|r_U(a)\)。
\end{proposition}
\begin{proof}
取两个归一化映射 \(f_U,f_V\)，则 \(f_V\circ f_U^{-1}\) 把圆盘映入圆盘并固定0。Schwarz 导数估计给 \(f_V'(a)/f_U'(a)\le1\)；取等时该复合是旋转，从而 \(f_V(U)=\mathbb D=f_V(V)\)，单射性迫使 \(U=V\)。共形变换公式由逆映射的链式法则直接得到。
\end{proof}
若 \(D(a,\rho)\subset U\subset D(a,R)\)，便有 \(\rho\le r_U(a)\le R\)。这个简单比较可在不求显式映射的情况下估计归一化导数。后面的 Poincaré 密度将等于 \(2/r_U(a)\)；本章先给出的极值解释会使度量的单调性自然出现。

共形半径还能度量区域对内点的影响。半径 \(R\) 的圆盘中，中心点共形半径为 \(R\)，靠近边界点时却趋于0；所以它描述的是“以当前点为中心的共形尺度”，并非一个区域全局共有的半径。对右半平面，向右平移一点使共形半径线性增大；沿竖直方向移动则保持不变，因为竖直平移是半平面的自同构。这些变化可由显式式子检验，也可由共形协变公式解释。

\ProseParagraphBreak
若区域增加了一条很远的狭窄通道，固定点的共形半径仍可能增加，但仅由集合包含关系只能断言严格单调，不能从图形面积的增加量读出改变量。共形几何的极值量关注全纯函数可以如何变化，而非欧氏面积。在同样面积的两个区域中，共形半径可以不同；在一个无界区域中它也可有限。这正是选取圆盘映射归一化导数作为不变量的意义。

\subsection{圆盘自同构}
\label{b3:subsec:050303}

\begin{proposition}[圆盘自同构]\label{b3:prop:05-disk-automorphisms}
每个圆盘自同构都唯一写成
\[
 h(z)=e^{\mathrm{i}\theta}\frac{z-a}{1-\overline a z},\qquad a\in\mathbb D,
\]
其中 \(e^{\mathrm{i}\theta}\) 是单位复数。反之，这些映射都是圆盘自同构。
\end{proposition}
\begin{proof}
取 \(a=h^{-1}(0)\)，则 \(h\circ\phi_a^{-1}\) 是固定0的圆盘自同构，由刚证明的 Schwarz 等号论证为旋转。反向由第一章的圆盘 Möbius 计算。参数 \(a\) 是唯一零点，旋转因子再由任一点或导数确定。
\end{proof}
圆盘自同构有三个实参数，而 \(f(a)=0\) 给两个实条件，导数正实给一个相位条件。这个参数解释与唯一性证明相符，但参数计数本身不是存在性或唯一性的证明。

\begin{theorem}\label{b3:thm:05-entire-univalent}
整个平面上的单叶整函数恰为 \(f(z)=az+b\)，其中 \(a\ne0\)。
\end{theorem}
\begin{proof}
若无穷远是本性奇点，Casorati--Weierstrass 定理给 \(|z_n|\rightarrow\infty\) 且 \(f(z_n)\rightarrow f(z_0)\)。在 \(z_0\) 附近单叶函数有全纯局部逆，因此充分晚的这些值同时有一个靠近 \(z_0\) 的原像。它与 \(z_n\) 不同，违反单射性。若无穷远可去，Liouville 定理使 \(f\) 常值，亦不可能。故无穷远为极点，\(f\) 是多项式。次数若至少为2，\(f'\) 由代数基本定理有零点，违反单叶函数导数处处非零。反向显然。
\end{proof}
于是全平面也不能双全纯映到任何真平面区域，即使目标无界。这比对有界区域使用 Liouville 更强。证明中的决定性限制是单射性；非常值整函数 \(e^z\) 的像是无界真区域 \(\mathbb C\setminus\{0\}\)，而它确实多对一。

共形类型的唯一性须把双全纯与覆盖区别开。圆盘可通过某些覆盖映到有孔平面域，但这并不说明两域双全纯。双全纯逆映射把任意路径收缩带回源域，因此保存单连通性；圆盘与去心圆盘不能双全纯。指数从平面覆盖去心平面也是多对一映射，与本章平面单叶函数分类不矛盾。第11章引入覆盖后，这种区分将成为统一化陈述的基本组成部分。

\subsection{共形分类}
\label{b3:subsec:050304}

真单连通平面区域因此只有一种共形类型，即圆盘类型。若把全平面也包括进来，平面与圆盘是不同类型，因为 Liouville 排除它们之间的双全纯映射。扩充球面又是第三种，它紧而平面与圆盘不紧；统一化理论将对一般单连通 Riemann 曲面说明这三种模型的穷尽性。

\ProseParagraphBreak

可计算的例子仍然值得保留。右半平面到圆盘可取 \((z-1)/(z+1)\)；水平带 \(|\operatorname{Im}z|<\pi/2\) 先由 \(e^z\) 单叶映到右半平面，再接前式；角域 \(0<\arg z<\alpha\) 先取明确分支 \(z^{\pi/\alpha}\) 映到上半平面，再用 Cayley 变换映到圆盘。每步都需要核对单射与满射，不能仅由边界线的像猜出整个区域的像。

\ProseParagraphBreak

多连通区域有真正的共形不变量。圆环不能映成圆盘，因为双全纯映射是同胚，保持闭曲线是否可收缩。\chapref{b3:ch:08}将以环域模量区分不同圆环；Riemann 映照定理的单连通条件正是在这一位置限制了分类范围。

裂平面的显式映射可以检验抽象证明中的平方根。对 \(U=\mathbb C\setminus(-\infty,0]\)，取 \(\operatorname{Re}\sqrt z>0\) 的平方根分支，再用右半平面的 Cayley 变换，得到
\[
 f(z)=\frac{\sqrt z-1}{\sqrt z+1}.
\]
它在1取0，导数为 \(1/4\)，所以 \(r_U(1)=4\)。点1到裂缝的欧氏距离为1，两种尺度相差显著。裂缝两侧会映到不同的圆周弧，即使它们在平面边界中对应同一负实点；第7章必须以素端处理这种边界辨识。

\ProseParagraphBreak
扇形 \(0<\arg z<\alpha\)、\(0<\alpha\le2\pi\)，可先用 \(z^{\pi/\alpha}=\exp((\pi/\alpha)\log z)\) 映到上半平面，再用分式变换映到圆盘。给定一个明确的辐角区间是必要数据。若忽略分支而直接把幂函数当成全平面函数，所得公式甚至可能不单值。对有理角度，代数式可用于计算；对任意角度，指数定义仍然统一有效。

Riemann 定理给存在性和唯一性，不给一般区域的闭式公式。对于圆盘、半平面、带和扇形，公式来自已知初等函数；对于不规则边界区域，极值构造仍然成立，却可能需要数值逼近才能观察映射。在判断一个候选公式时应逐项核对：在源域单值全纯、导数非零、整体单射、像确为整个目标，以及指定点处的归一化。只画出几个点的像或检查导数，不能替代满射与单射的证明。

\ProseParagraphBreak
例如 \(z\mapsto z^2\) 在右半平面单叶，像为去掉非正实轴的裂平面；若把源域换成上半平面，像则去掉非负实轴。公式相同而像域不同，是因为辐角区间改变。再把源域放宽到全平面，映射变成二对一并在0分歧。所有这些现象都可由辐角范围和局部零点阶数直接核对，体现了“一个函数公式”与“一张区域之间的共形等价”之间的差别。

所有这些显式映射都应在写出逆函数后再核对满射。带的指数逆取对数分支，扇形的幂逆也取指定辐角范围的根。逆的分支范围与正向映射的辐角范围恰好匹配，才保证每个目标点有且只有一个源点。
